\bookchapter{第七章：前沿功能全解析}

\chapterintro{科技日新月異，手機功能都喺度不斷進化。衞星通信、無網通、車禍檢測、星閃技術……呢啲聽落好科幻嘅功能，其實已經走入我哋嘅生活。等我哋一齊嚟瞭解呢啲前沿技術。}

\booksection{一、衞星通信：冇信號都可以求救}

衞星通信係近年嚟手機領域嘅一大突破。當你喺荒郊野外、深山老林，手機完全冇信號嗰陣，衞星通信可能就係救命稻草。

\booksubsection{工作原理：}

手機直接與軌道上嘅通信衞星建立連接，發送求救資訊。目前主要支援發送短訊同位置資訊，唔可以打電話。

\booksubsection{支援機型：}

\begin{itemize}
  \item \textbf{蘋果}：iPhone 14系列及以後，支援通過衞星發送緊急SOS。iPhone 17系列進一步支援衞星發送iMessage同短訊，且擴展咗更多地區嘅覆蓋。
  \item \textbf{華為}：Mate 50系列開始支援北斗衞星消息（僅文字），Mate 60 Pro系列進一步支援天通衞星通話（可打電話），Mate 70系列衞星通信能力進一步增強。
  \item \textbf{小米/OPPO/vivo}：2024年起旗艦機型陸續支援北斗衞星消息功能。
  \item \textbf{發展趨勢}：衞星通信正由高端旗艦專屬功能向下普及，未來中端機型都有望支援。中國北斗衞星系統喺國產手機上嘅支援最為廣泛。
\end{itemize}

\booksubsection{使用場景：}

\begin{itemize}
  \item 戶外探險、登山遇險。
  \item 海上遇險。
  \item 自然災害導致地面通信中斷。
\end{itemize}

\booksubsection{使用方法：}

\begin{enumerate}
  \item 確保喺開闊地帶，可以睇到天空。
  \item 打開手機嘅衞星通信功能（通常喺緊急呼叫介面）。
  \item 按照螢幕指引，對準衞星方向。
  \item 發送預設嘅求救資訊同位置。
\end{enumerate}

\begin{tcbinfo}{小知識}
衞星通信需要手機有特殊嘅硬件天線，唔係所有手機都支援。另外，衞星通信通常需要付費訂閲服務，但緊急求救功能一般係免費嘅。

\end{tcbinfo}

\booksection{二、無網通：近距離通信新方式}

「無網通」係指喺冇流動網絡同Wi-Fi嘅情況下，兩台手機之間直接通信嘅技術。

\booksubsection{技術原理：}

利用藍牙或專用嘅無線電頻率，喺短距離內（通常幾百米到幾公里）建立點對點連接。

\booksubsection{應用場景：}

\begin{itemize}
  \item \textbf{戶外活動}：徒步、露營嗰陣隊友間通信。
  \item \textbf{災難救援}：地震、洪水後通信中斷嗰陣。
  \item \textbf{大型活動}：演唱會、體育賽事人多信號差嗰陣。
\end{itemize}

\booksubsection{支援平臺：}

\begin{itemize}
  \item \textbf{蘋果}：iOS 17嘅「NameDrop」同「離線揾朋友」功能。
  \item \textbf{華為}：暢連App支援無網通話。
  \item \textbf{小米}：小米妙享中心嘅部分功能。
\end{itemize}

\booksection{三、車禍檢測：關鍵時刻救命}

車禍檢測係利用手機傳感器檢測嚴重車禍，並自動撥打緊急電話嘅功能。

\booksubsection{工作原理：}

手機通過加速度計、陀螺儀、氣壓計、麥克風等多個傳感器，綜合判斷係咪發生咗嚴重碰撞。檢測到車禍後，手機會：

\begin{enumerate}
  \item 彈出倒計時介面，詢問你係咪冇事。
  \item 如果倒計時結束你冇回應，自動撥打緊急電話。
  \item 向緊急聯絡人發送你嘅位置。
\end{enumerate}

\booksubsection{支援機型：}

\begin{itemize}
  \item \textbf{蘋果}：iPhone 14系列及以後。
  \item \textbf{安卓}：部分高端機型陸續支援。
\end{itemize}

\booksubsection{注意事項：}

\begin{itemize}
  \item 功能需要開啟先可以使用。
  \item 可能會誤觸發（比如過山車、劇烈運動），及時取消即可。
  \item 依賴傳感器同算法，唔係100\%可靠，但多一層保障。
\end{itemize}

\booksection{四、星閃技術：下一代無線連接}

星閃（SparkLink）係華為牽頭研發嘅新一代無線通信技術，目標係喺某些場景下替代藍牙同Wi-Fi。佢喺關鍵指標上對標甚至超越現有方案：空口速率可達900Mbps（遠超藍牙嘅50Mbps），延遲低至20 $\mu$ s（藍牙約15ms），單網絡支援4096個設備同時連接（藍牙通常8個以內）。

\booksubsection{技術特點：}

\begin{itemize}
  \item \textbf{超低延遲}：端到端延遲僅20微秒（0.02ms），比藍牙低兩個數量級，遊戲同音樂同步幾乎無感知。
  \item \textbf{高速率}：單鏈路峯值速率900Mbps，係藍牙音頻嘅數十倍，可傳輸無損音頻甚至高清影片。
  \item \textbf{低功耗}：待機功耗與藍牙相當，但傳輸效率更高，單位數據嘅能耗更低。
  \item \textbf{廣覆蓋+高並發}：連接距離比藍牙遠50\%以上，單網絡支援上千設備同時在線。
\end{itemize}

\booksubsection{應用場景：}

\begin{itemize}
  \item \textbf{無線耳機}：更低延遲，玩遊戲音畫同步。
  \item \textbf{智能手錶}：更穩定嘅連接。
  \item \textbf{智能家居}：更多設備同時連接。
  \item \textbf{車載互聯}：更流暢嘅手機-車機連接。
\end{itemize}

\booksubsection{支援設備：}

目前主要係華為生態設備，如華為Mate 60系列、部分華為耳機同手錶。未來會有更多品牌支援。

\textbf{技術展望}\par

星閃技術代表咗無線通信嘅未來方向。隨住更多廠商加入，我哋可能會見到更多支援星閃嘅設備，享受更好嘅無線體驗。

\begin{tcbpractice}{去檢查一下：你嘅手機係咪支援呢啲功能}
\begin{enumerate}
  \item 去設定入面搜尋「衞星」，睇下有冇衞星通信相關選項。
  \item 搜尋「車禍檢測」或「緊急SOS」，檢查係咪開啟。
  \item 如果你嘅手機係華為，睇下係咪支援星閃設備連接。
  \item 瞭解呢啲功能嘅使用方法，關鍵時刻可能派上用場。
\end{enumerate}

\end{tcbpractice}
